\documentclass{article}
\usepackage{amsmath}
\usepackage{amssymb}
\usepackage{mathtools}
\usepackage{graphicx} % Required for inserting images
\usepackage[margin=1in]{geometry}
\usepackage{float}
\begin{document}

\section{Variables}

\subsection{Constants}

\begin{equation}
\begin{aligned}
     \Vec{D}_{fl} \coloneqq& \text{Displacement from center-of-mass of front-left rotor (3-dimensional)}\\
     \Vec{N} \coloneqq& \text{Rotor plane normal vector (fixed to the +z unit vector)}
\end{aligned}
\end{equation}

\subsection{Inputs}
\begin{equation}
\begin{aligned}
     V_{fl} \coloneqq& \text{Voltage applied to front left motor}
\end{aligned}
\end{equation}

\subsection{Outputs}
\begin{equation}
\begin{aligned}
     \omega_{fl} \coloneqq& \text{Angular velocity of front left rotor}\\
     \dot{\omega_{fl}} \coloneqq& \text{Angular acceleration of front left rotor}\\
     \vec{T}_{fl} \coloneqq& \text{Thrust of the front left rotor}\\
     \vec{\tau}_{fl} \coloneqq& \text{Torque vector of the front left rotor}\\
     \gamma \coloneqq& \text{Yaw angle of drone}\\
     \phi \coloneqq& \text{Pitch angle of drone}\\
     \theta \coloneqq& \text{Roll angle of drone}\\
\end{aligned}
\end{equation}

\section{Dynamics}

\subsection{Rotors}

\begin{equation}
\begin{aligned}
    V =& k_{\omega}(\omega) + k_{\omega^2}(\omega^2) + k_{\alpha}(\dot{\omega})\\
    \vec{T} =& k_T(\omega_m^2)\vec{N}\\
    \vec{\tau} =& k_{\tau}(\omega_m^2)\vec{N}\\
\end{aligned}
\end{equation}

\subsection{Airframe}

\subsubsection{Yaw}
\begin{equation}
\begin{aligned}
    \ddot{\gamma} =& \rho_{yaw,lin}(\dot{\gamma}) + \rho_{yaw,quad}(\dot{\gamma}^2) + \sum_{m \in fl, fr, bl, br} \pm \left(k_{inertia} (\dot{\omega_{m}}) + k_\tau(\omega_{m}^2)\right)
\end{aligned}
\end{equation}
where the sign represents whether sensor-forward for the rotor is clockwise or anticlockwise (we adopt the convention that positive RPM corresponds to positive thrust, so this depends on the diagonal on which the rotor lies).

Regrouping by terms, we obtain the regression equation:

\begin{equation}
\begin{aligned}
    \ddot{\gamma} = \rho_{yaw,lin}(\dot{\gamma}) + \rho_{yaw,quad}(\dot{\gamma}^2) +& k_{inertia}(\dot{\omega_{fl}} - \dot{\omega_{fr}} - \dot{\omega_{bl}} + \dot{\omega_{br}})\\
    +& k_\tau(\omega_{fl}^2 - \omega_{fr}^2 - \omega_{bl}^2 + \omega_{br}^2)
\end{aligned}
\end{equation}

\subsubsection{Pitch and Roll}
Signs here are given for pitch by front (+) vs back (-) and for roll by left (+) vs right (-). 
\begin{equation}
\begin{aligned}
    \ddot{\phi} =& \rho_{pitch,lin}(\dot{\phi}) + \rho_{pitch,quad}(\dot{\phi}^2) + \sum_{m \in fl, fr, bl, br} \pm ||\vec{D_m} \times \vec{x}||\left(k_T(\omega_{m}^2)\right)\\
    \ddot{\theta} =& \rho_{roll,lin}(\dot{\theta}) + \rho_{roll,quad}(\dot{\theta}^2) + \sum_{m \in fl, fr, bl, br} \pm ||\vec{D_m} \times \vec{y}||\left(k_T(\omega_{m}^2)\right)
\end{aligned}
\end{equation}

\subsection{Linearization}
By Taylor expansion:
\begin{equation}
\begin{aligned}
    ax + bx^2 &=_{x \sim x_0} (bx_0^2 + 2bx_0(x-x_0)) + ax\\
    &= (a + 2bx_0)x - bx_0^2
\end{aligned}
\end{equation}
We expand at $\omega_m \sim \widehat{\omega}_m$, and for compactness we write

\begin{align*}
    \rho_y =& \rho_{yaw,lin}\\
    \rho_p =& \rho_{pitch,lin}\\
    \rho_r =& \rho_{roll,lin}\\
    \{T, \tau\}_{1, 2} =& \{T, \tau\}_{lin, quad}\\
    k_I =& k_{inertial} \\
    k_{\beta, m} =& -\frac{1}{k_\alpha} \left(k_\omega + 2 k_{\omega^2} \widehat{\omega}_m \right)
\end{align*}
where $k_{\beta, m}$ is justified by the rewriting of the voltage equation on the rotors as
\begin{equation}
    \begin{aligned}
     \dot{\omega}_m &= \frac{1}{k_\alpha} \left(V - k_\omega(\omega_m) - k_{\omega^2}(\omega^2)\right)\\
     &=_{\omega_m \sim \widehat{\omega}_m} \frac{1}{k_\alpha} \left( V - \left(k_\omega + 2 k_{\omega^2}\widehat{\omega}_m\right) \omega_m + k_{\omega^2}\widehat{\omega}_m^2 \right)\\
     &= - \frac{1}{k_\alpha} \left(k_\omega + 2 k_{\omega^2}\widehat{\omega}_m\right) \omega_m + \frac{1}{k_\alpha} \left( V + k_{\omega^2}\widehat{\omega}_m^2 \right),
    \end{aligned}
\end{equation}
that is, our $k_{\beta, m}$ is the coefficient of $\omega$, where the Taylor expansion is given by the above linearization; the voltage-dependent term and the constant term also appear in the appropriate matrices.

Additionally, we assume some basic symmetry in the prop arrangement and adopt the following simplification:

\begin{align*}
    D_F \coloneqq& ||\vec{D_{FL}} \times \vec{y}|| = ||\vec{D_{FR}} \times \vec{y}||\\
    D_B \coloneqq& ||\vec{D_{BL}} \times \vec{y}|| = ||\vec{D_{BR}} \times \vec{y}||\\
    D_L \coloneqq& ||\vec{D_{FL}} \times \vec{x}|| = ||\vec{D_{BL}} \times \vec{x}||\\
    D_R \coloneqq& ||\vec{D_{FR}} \times \vec{x}|| = ||\vec{D_{BR}} \times \vec{x}||\\
\end{align*}

We thus define our state vector

\setcounter{MaxMatrixCols}{20}
\begin{equation}
\begin{aligned}
    X \coloneqq \begin{bmatrix}
    \dot{\gamma} & \dot{\phi} & \dot{\theta} & \omega_{FL} & \omega_{FR} & \omega_{BL} & \omega_{BR}
\end{bmatrix},
\end{aligned}
\end{equation}

and transition matrix

\begin{equation}
\hspace{-0.7in}A \coloneqq \begin{bmatrix}
\rho_y & 0 & 0 & 2 k_\tau \widehat{\omega}_{FL} + k_Ik_{\beta, FL} & - 2 k_\tau \widehat{\omega}_{FR} - k_Ik_{\beta, FR} & - 2 k_\tau \widehat{\omega}_{BL} - k_Ik_{\beta, BL} & 2 k_\tau \widehat{\omega}_{BR} + k_Ik_{\beta, BR}\\
0 & \rho_p & 0 & 2 k_T D_F \hat{\omega_{FL}} & 2 k_T D_F \widehat{\omega}_{FR} & -2 k_T D_B \widehat{\omega}_{BL} & -2 k_T D_B\widehat{\omega}_{BR}\\
0 & 0 & \rho_r & 2 k_T D_L\hat{\omega_{FL}} & -2 k_T D_R \widehat{\omega}_{FR} & 2 k_T D_L\widehat{\omega}_{BL} & -2 k_T D_R\widehat{\omega}_{BR}\\
0 & 0 & 0 & k_{\beta, FL} & 0 & 0 & 0\\
0 & 0 & 0 & 0 & k_{\beta, FR} & 0 & 0 &\\
0 & 0 & 0 & 0 & 0 & k_{\beta, BL} & 0\\
0 & 0 & 0 & 0 & 0 & 0 & k_{\beta, BR}\\
\end{bmatrix}.\\
\end{equation}

Our Taylor expansion also leaves an affine vector term of the form

\begin{equation}
\begin{aligned} C \coloneqq 
    \begin{bmatrix}
        & \left(\dfrac{k_I k_{\omega^2}}{k_\alpha} - k_\tau\right) \left( \widehat{\omega}_{FL}^2 - \widehat{\omega}_{FR}^2 - \widehat{\omega}_{BL}^2 + \widehat{\omega}_{BR}^2\right)\\
        & -k_T D_F(\widehat{\omega_{FL}}^2 + \widehat{\omega_{FR}}^2) + k_T D_B(\widehat{\omega_{BL}}^2 + \widehat{\omega_{BR}}^2)\\
        & - k_T D_L(\widehat{\omega_{FL}}^2 + \widehat{\omega_{BL}}^2) + k_T D_R(\widehat{\omega_{FR}}^2 + \widehat{\omega_{BR}}^2)\\
        & \dfrac{k_{\omega^2}\widehat{\omega}_{FL}^2}{k_\alpha}\\
        & \dfrac{k_{\omega^2}\widehat{\omega}_{FR}^2}{k_\alpha}\\
        & \dfrac{k_{\omega^2}\widehat{\omega}_{BL}^2}{k_\alpha}\\
        & \dfrac{k_{\omega^2}\widehat{\omega}_{BR}^2}{k_\alpha}\\
    \end{bmatrix}.
\end{aligned} 
\end{equation}

Lastly, we define the control vector, given by

\begin{equation}
\begin{aligned}U \coloneqq 
    \begin{bmatrix}
        V_{FL} & V_{FR} & V_{BL} & V_{BR}
    \end{bmatrix}
\end{aligned} 
\end{equation}

with input matrix

\begin{equation}
\begin{aligned}B \coloneqq 
    \begin{bmatrix}
        \dfrac{k_I}{k_\alpha}  & -\dfrac{k_I}{k_\alpha} & -\dfrac{k_I}{k_\alpha} & \dfrac{k_I}{k_\alpha} \\
        0 & 0 & 0 & 0\\
        0 & 0 & 0 & 0\\
        \frac{1}{k_\alpha} & 0 & 0 & 0\\
        0 & \frac{1}{k_\alpha} & 0 & 0\\
        0 & 0 & \frac{1}{k_\alpha} & 0\\
        0 & 0 & 0 & \frac{1}{k_\alpha}\\
    \end{bmatrix}.
\end{aligned} 
\end{equation}\\


Together, the above yield our continuous-time state-space equation:
\begin{equation}
    \dot{X} = AX + BU + C,
\end{equation}
where we treat $X, U$ as column vectors as appropriate.

\subsection{Closed-loop linearization}
The affine term is cancelled by the input
\begin{equation}
    U_{aff} \coloneqq -B^+ C(\widehat{\omega}),
\end{equation}
which removes the component of $C$ in the range of $B$ (the yaw and motor rows; $B$ has no direct pitch or roll path, so those rows of $C$ are left to the integrator).  Were $\widehat{\omega}$ a fixed operating point, $U_{aff}$ would be a constant offset and the dynamics seen by the feedback would be exactly $\dot{X} = AX + BU$.  In practice, however, we evaluate the linearization at the \emph{measured} rotor speeds on every control tick, $\widehat{\omega} \coloneqq \omega$, so $U_{aff}$ is itself a function of the state, i.e.\ feedback, and it contributes to the closed-loop Jacobian.

For a single rotor, the true dynamics and the linearizer give
\begin{equation}
\begin{aligned}
    k_\alpha \dot{\omega}_m &= V_{lin} + V_{aff} - k_\omega(\omega_m) - k_{\omega^2}(\omega_m^2)\\
    &= V_{lin} - k_{\omega^2}(\omega_m^2) - k_\omega(\omega_m) - k_{\omega^2}(\omega_m^2)\\
    &= V_{lin} - k_\omega(\omega_m) - 2k_{\omega^2}(\omega_m^2),
\end{aligned}
\end{equation}
which agrees with the tangent line $k_\alpha(k_{\beta,m}\,\omega_m) + V_{lin}$ in value at $\omega_m = \widehat{\omega}_m$ but has twice its curvature, and hence slope
\begin{equation}
    \frac{\partial \dot{\omega}_m}{\partial \omega_m} = -\frac{1}{k_\alpha}\left(k_\omega + 4k_{\omega^2}\widehat{\omega}_m\right) \neq k_{\beta, m}.
\end{equation}
The discrepancy is first-order in $\delta\omega_m = \omega_m - \widehat{\omega}_m$: expanding the intercept $c(\widehat{\omega}_m) = k_{\omega^2}\widehat{\omega}_m^2 / k_\alpha$ (the motor rows of $C$),
\begin{equation}
    c(\widehat{\omega}_m + \delta\omega_m) = c(\widehat{\omega}_m) + \underbrace{\frac{2k_{\omega^2}\widehat{\omega}_m}{k_\alpha}\delta\omega_m}_{\text{first order}} + \underbrace{\frac{k_{\omega^2}}{k_\alpha}\delta\omega_m^2}_{\text{dropped}},
\end{equation}
and cancelling $c$ at the moving point cancels its linear part as well.  Note that because the tangent line passes through the true dynamics at $\widehat{\omega}_m$,
\begin{equation}
    \frac{d}{d\widehat{\omega}_m}\left[k_{\beta,m}(\widehat{\omega}_m)\,\widehat{\omega}_m + c(\widehat{\omega}_m)\right] = k_{\beta, m}
    \quad\Longrightarrow\quad
    k_{\beta,m}'(\widehat{\omega}_m)\,\widehat{\omega}_m + c'(\widehat{\omega}_m) = 0,
\end{equation}
so for the plant alone the moving slope and moving intercept cancel, and $A$ is its true Jacobian; it is removing only the intercept that leaves the residual $-c'$.

In matrix form, the closed loop is $\dot{X} = f(X) + BU_{lin} - BB^+C(\omega)$, with Jacobian
\begin{equation}
    A_{eff} \coloneqq A - BB^+\frac{\partial C}{\partial \omega}\bigg|_{\widehat{\omega}}
    = A + B\,G\,\mathrm{diag}\left(2|\widehat{\omega}_m|\right),
    \qquad G \coloneqq -B^+ P,\quad C = P\,[\widehat{\omega}_m|\widehat{\omega}_m|].
\end{equation}
$BB^+$ leaves the pitch and roll rows untouched, passes the motor rows, and carries them into the yaw row through $k_I$.  To within the least-squares share of the $-k_\tau$ yaw residual (of order $10^{-5}$ relative), $G = -k_{\omega^2}I$, and the correction amounts to the substitution
\begin{equation}
    k_{\beta, m} \;\to\; \tilde{k}_{\beta, m} \coloneqq -\frac{1}{k_\alpha}\left(k_\omega + 4 k_{\omega^2}\widehat{\omega}_m\right)
\end{equation}
wherever $k_{\beta, m}$ appears in $A$ (the motor diagonal and the $k_I k_{\beta, m}$ terms of the yaw row).  The gains, the integral bound and the feedforward reference model are all designed on $A_{eff}$.


\end{document}
